% Tradução brasileira revista a partir do workpass espanhol congelado;
% a fonte permanece em source_es.
% SGA 3, Exposé VIII, tradução brasileira: §4 até a primeira frase
% da demonstração do Corolário 4.7.
% Autoridade: Polo--Gille Exp8-8nov09.pdf, pp. locais 11--13,
% leitor completo pp. 627--629.

\SGARefTarget{sga3:viii:section:4}\section*{4. Torsores sob um grupo diagonalizável}
\addcontentsline{toc}{section}{4. Torsores sob um grupo diagonalizável}

\SGARefTarget{sga3:viii:section:4:problem:01}Seja \(S\) um pré-esquema e seja
\(G=D_S(M)\) um grupo diagonalizável sobre \(S\). Determinaremos os
\(G\)-torsores (ou fibrados principais homogêneos sob \(G\)) sobre \(S\)
para a topologia fielmente plana e quase compacta; cf. Exposé IV,
\SGARefLink{sga3:iv:subsection:5:1}{5.1}. Recordemos que um pré-esquema
\(P\) sobre \(S\), com grupo de operadores \(S\)-grupal \(G\), chama-se
\emph{torsor} ou \emph{fibrado principal homogêneo} se todo ponto de \(S\)
admite uma vizinhança aberta \(U\) e um morfismo fielmente plano e
quase compacto \(S'\to U\) tais que
\[
  P'=P\times_SS'
\]
é isomorfo, como objeto com operadores, a
\[
  G'=G\times_SS',
\]
onde \(G'\) atua sobre si mesmo por translações à direita. Como
\(G\) é afim sobre \(S\), SGA 1, VIII 5.6 implica que \(P\) é
necessariamente afim sobre \(S\). Além disso, como o próprio \(G\) é fielmente
plano e quase compacto sobre \(S\), \(P\) é principal homogêneo sob \(G\)
se, e somente se, for formalmente principal homogêneo e, além disso, fielmente plano e
quase compacto sobre \(S\) (cf. IV,
\SGARefLink{sga3:iv:proposition:5:1:6}{5.1.6}).

Recordemos também
(\SGARefLink{sga3:i:proposition:4-7-3}{Exposé I, 4.7.3}) que dar um
\(S\)-pré-esquema \(P\), afim sobre \(S\), com grupo de operadores
\(G=D_S(M)\), equivale a dar uma álgebra comutativa quase-coerente
graduada por \(M\) sobre \(S\). Em outras palavras, equivale a dar uma
álgebra quase-coerente \(\mathscr A\) sobre \(S\), juntamente com uma
decomposição em soma direta como módulo
\[
  \mathscr A=\coprod_{m\in M}\mathscr A_m,
\]
tal que
\[
  \mathscr A_m\mathbin{\cdot}\mathscr A_{m'}
  \subseteq\mathscr A_{m+m'}
  \qquad\text{para }m,m'\in M.
\]
Feita essa observação, a resposta ao
\SGARefLink{sga3:viii:section:4:problem:01}{problema anterior} é a
seguinte.

\medskip
\noindent\SGARefTarget{sga3:viii:proposition:4:1}\textbf{Proposição 4.1.}
\textit{Seja \(P\) o pré-esquema com grupo de operadores sobre \(S\)
\(G=D_S(M)\), definido pela álgebra \(M\)-graduada \(\mathscr A\).
Então \(P\) é um fibrado principal homogêneo sob \(G\) se, e somente se,
\(\mathscr A\) satisfaz as condições seguintes:}
\begin{enumerate}
  \item[(a)]\SGARefTarget{sga3:viii:proposition:4:1:item:a:01}
  \textit{para todo \(m\in M\), \(\mathscr A_m\) é um módulo invertível
  sobre \(S\);}

  \item[(b)]\SGARefTarget{sga3:viii:proposition:4:1:item:b:01}
  \textit{para \(m,m'\in M\), o homomorfismo}
  \[
    \mathscr A_m\otimes_{\mathcal O_S}\mathscr A_{m'}
    \longrightarrow\mathscr A_{m+m'}
  \]
  \textit{induzido pela multiplicação em \(\mathscr A\) é um
  isomorfismo.}
\end{enumerate}

A necessidade de \SGARefLink{sga3:viii:proposition:4:1}{essas condições}
deduz-se imediatamente por descida, pois elas são satisfeitas quando \(P\)
é o fibrado principal homogêneo trivial, isto é, quando
\(\mathscr A=\mathcal O_S(M)\). Para a suficiência, observe-se que
\SGARefLink{sga3:viii:proposition:4:1:item:a:01}{(a)} já implica que \(P\)
é fielmente plano sobre \(S\); de qualquer modo, é quase compacto sobre \(S\),
pois é afim sobre \(S\). Resta, portanto, verificar que \(P\) é
formalmente principal homogêneo sob \(G\), isto é, que o morfismo usual
\SGARefTarget{sga3:viii:proposition:4:1:diagram:01}
\[
  P\times_SG\longrightarrow P\times_SP
\]
é um isomorfismo. Em termos de álgebras graduadas,
\SGARefLink{sga3:viii:proposition:4:1:diagram:01}{esse morfismo} é dado
explicitamente pelo homomorfismo
\SGARefTarget{sga3:viii:proposition:4:1:diagram:02}
\[
  \mathscr A\otimes\mathscr A
  \longrightarrow
  \mathscr A(M)=\mathscr A\otimes\mathcal O_S(M).
\]

No bigrau \((m,n)\), onde \(m,n\in M\), ele é dado por
\[
  x_m\otimes y_n\longmapsto x_my_n\otimes e_n.
\]
No nível dos graus,
\SGARefLink{sga3:viii:proposition:4:1:diagram:02}{esse homomorfismo} é
compatível com o homomorfismo \(M\times M\to M\times M\) dado por
\[
  (m,n)\longmapsto(m+n,n),
\]
o que mostra que
\SGARefLink{sga3:viii:proposition:4:1:item:b:01}{(b)} exprime precisamente,
e independentemente de
\SGARefLink{sga3:viii:proposition:4:1:item:a:01}{(a)}, que \(P\) é
formalmente principal homogêneo, e prova
\SGARefLink{sga3:viii:proposition:4:1}{4.1}.

A descida fielmente plana também fornece:

\medskip
\noindent\SGARefTarget{sga3:viii:corollary:4:2}\textbf{Corolário 4.2.}
\textit{As condições de \SGARefLink{sga3:viii:proposition:4:1}{4.1}
implicam que o homomorfismo}
\[
  \mathcal O_S\longrightarrow\mathscr A_0
\]
\textit{é um isomorfismo.}

Por exemplo, se \(M=\mathbb Z\), então, sob as condições de
\SGARefLink{sga3:viii:proposition:4:1}{4.1}, a álgebra \(\mathscr A\) fica
essencialmente determinada por \(\mathscr A_1=\mathscr L\), a saber,
\[
  \mathscr A\simeq\coprod_{n\in\mathbb Z}\mathscr L^{\otimes n}
\]
como álgebras graduadas. Obtemos assim o resultado conhecido:

\medskip
\noindent\SGARefTarget{sga3:viii:corollary:4:3}\textbf{Corolário 4.3.}
\textit{A categoria dos fibrados principais homogêneos \(P\) sobre
\(S\) sob \(\mathbb G_{m,S}\) é equivalente à categoria dos
módulos invertíveis \(\mathscr L\) sobre \(S\), sendo os morfismos em
ambas as categorias os isomorfismos das estruturas consideradas. Obtêm-se
funtores quase inversos associando a \(P\) a componente de grau
\(1\) de sua álgebra afim \(\mathbb Z\)-graduada, e associando a
\(\mathscr L\) o espectro da álgebra \(\mathbb Z\)-graduada}
\[
  \coprod_{n\in\mathbb Z}\mathscr L^{\otimes n}.
\]

Em particular:

\medskip
\noindent\SGARefTarget{sga3:viii:corollary:4:4}\textbf{Corolário 4.4.}
\textit{O grupo das classes de isomorfia de fibrados principais homogêneos
sobre \(S\) sob \(\mathbb G_{m,S}\) é isomorfo ao grupo
\(\operatorname{Pic}(S)\) das classes de isomorfia de módulos invertíveis
sobre \(S\), isto é, a}
\[
  H^1(S,\mathcal O_S^\times).
\]

Como \(\mathbb G_{m,S}\) é o esquema de automorfismos do módulo
\(\mathcal O_S\), o Corolário
\SGARefLink{sga3:viii:corollary:4:4}{4.4} equivale ao
\SGARefLink{sga3:viii:corollary:4:5}{enunciado seguinte}, uma das formas
do «\SGARefLink{sga3:viii:corollary:4:5}{Teorema 90}» de Hilbert.

\medskip
\noindent\SGARefTarget{sga3:viii:corollary:4:5}\textbf{Corolário 4.5.}
\textit{Todo fibrado principal homogêneo sobre \(S\) sob \(\mathbb G_m\)
é localmente trivial para a topologia de Zariski.}

\medskip
\noindent\SGARefTarget{sga3:viii:remark:4:5:1}\textbf{Observação 4.5.1.}
\SGARefLink{sga3:viii:corollary:4:5}{O enunciado anterior} deixa de ser
verdadeiro, em geral, para um grupo como \(\mu_n\), ou para uma «forma
torcida» de \(\mathbb G_m\). Por exemplo, a única forma torcida de
\(\mathbb G_m\) sobre o corpo \(\mathbb R\) dos números reais tem
primeiro grupo de coomologia igual a \(\mathbb Z/2\mathbb Z\).

\footnote{Nota do editor: foi acrescentado o que segue.}
De fato, seja \(\mathbb S^1\) o núcleo do morfismo norma
\[
  N:\prod_{\mathbb C/\mathbb R}\mathbb G_{m,\mathbb C}
  \longrightarrow\mathbb G_{m,\mathbb R};
\]
é uma forma torcida por \(\mathbb C/\mathbb R\) de
\(\mathbb G_{m,\mathbb R}\). A equação \(N(z)=-1\) em
\(\prod_{\mathbb C/\mathbb R}(\mathbb G_{m,\mathbb C})\) define um torsor
\(X\) sob \(\mathbb S^1\) sobre \(\operatorname{Spec}(\mathbb R)\),
localmente trivial para a topologia étale, mas não trivial, pois
\(X(\mathbb R)=\varnothing\). Mostremos que
\[
  H^1_{\mathrm{\acute et}}(\mathbb R,\mathbb S^1)
  \simeq\mathbb Z/2\mathbb Z.
\]
Há uma sequência exata de grupos algébricos comutativos lisos sobre
\(\mathbb R\)
\SGARefTarget{sga3:viii:remark:4:5:1:diagram:01}
\[
  1\longrightarrow\mathbb S^1
  \longrightarrow\prod_{\mathbb C/\mathbb R}\mathbb G_{m,\mathbb C}
  \longrightarrow\mathbb G_{m,\mathbb R}
  \longrightarrow1,
\]
que dá origem a uma sequência exata longa de coomologia étale (ou de
coomologia de Galois):
\SGARefTarget{sga3:viii:remark:4:5:1:diagram:02}
\[
\begin{aligned}
  0\longrightarrow\mathbb S^1(\mathbb R)
  &\longrightarrow\mathbb C^\times
   \xrightarrow{N}\mathbb R^\times
   \longrightarrow H^1_{\mathrm{\acute et}}(\mathbb R,\mathbb S^1)\\
  &\longrightarrow
   H^1_{\mathrm{\acute et}}\!
   \left(\mathbb R,
     \prod_{\mathbb C/\mathbb R}\mathbb G_{m,\mathbb C}\right)
   \longrightarrow\cdots .
\end{aligned}
\]
Ora (veja-se, por exemplo, XXIV,
\SGARefLink{sga3:xxiv:proposition:8-4}{8.4}),
\[
  H^1_{\mathrm{\acute et}}\!
  \left(\mathbb R,
    \prod_{\mathbb C/\mathbb R}\mathbb G_{m,\mathbb C}\right)
  \simeq
  H^1_{\mathrm{\acute et}}(\mathbb C,\mathbb G_{m,\mathbb C}),
\]
e este último grupo é zero por
\SGARefLink{sga3:viii:corollary:4:5}{4.5} (ou, neste caso, porque
\(\mathbb C\) é algebricamente fechado). Obtém-se, portanto, um
isomorfismo
\[
  H^1_{\mathrm{\acute et}}(\mathbb R,\mathbb S^1)
  \simeq\mathbb R^\times/N(\mathbb C^\times)
  \simeq\{\pm1\}.
\]

Precisaremos do \SGARefLink{sga3:viii:proposition:4:6}{resultado seguinte}
na \SGARefLink{sga3:viii:section:5}{seção seguinte}.

\medskip
\noindent\SGARefTarget{sga3:viii:proposition:4:6}\textbf{Proposição 4.6.}
\textit{Sob as condições de
\SGARefLink{sga3:viii:proposition:4:1}{4.1}, as condições
\SGARefLink{sga3:viii:proposition:4:1:item:a:01}{(a)} e
\SGARefLink{sga3:viii:proposition:4:1:item:b:01}{(b)} equivalem às
condições seguintes:}
\begin{enumerate}
  \item[\((a')\)]\SGARefTarget{sga3:viii:proposition:4:6:item:aprime:01}
  \textit{o homomorfismo \(\mathcal O_S\to\mathscr A_0\) é um
  isomorfismo;}

  \item[\((b')\)]\SGARefTarget{sga3:viii:proposition:4:6:item:bprime:01}
  \textit{para todo \(m\in M\) (basta exigi-lo para os elementos de um
  sistema de geradores de \(M\)), tem-se}
  \[
    \mathscr A_m\mathbin{\cdot}\mathscr A_{-m}=\mathscr A_0.
  \]
\end{enumerate}

A necessidade é clara em virtude de
\SGARefLink{sga3:viii:corollary:4:2}{4.2}.\footnote{Nota do editor: foi
corrigido «\SGARefLink{sga3:viii:proposition:4:1}{4.1}» para
«\SGARefLink{sga3:viii:corollary:4:2}{4.2}».} Ficamos assim reduzidos a
demonstrar o \SGARefLink{sga3:viii:corollary:4:7}{corolário seguinte}.

\medskip
\noindent\SGARefTarget{sga3:viii:corollary:4:7}\textbf{Corolário 4.7.}
\textit{Seja}
\[
  A=\coprod_{n\in\mathbb Z}A_n
\]
\textit{um anel \(\mathbb Z\)-graduado tal que}
\[
  A_1\mathbin{\cdot}A_{-1}=A_0.
\]
\textit{Então os \(A_n\) são \(A_0\)-módulos invertíveis e, para
\(n,n'\in\mathbb Z\), o homomorfismo}
\[
  A_n\otimes_{A_0}A_{n'}\longrightarrow A_{n+n'}
\]
\textit{induzido pela multiplicação em \(A\) é um isomorfismo.}

Por hipótese, existem \(f_i\in A_1\) e \(g_i\in A_{-1}\) tais que
\SGARefTarget{sga3:viii:corollary:4:7:eq:star:01}
\[
  \makebox[\displaywidth][s]{%
    \(\text{(*)}\)\hfill
    \(\displaystyle\sum_i f_ig_i=1.\)\hfill
    \(\phantom{\text{(*)}}\)%
  }
\]
Como a afirmação a demonstrar é local sobre
\(\operatorname{Spec}(A_0)\),\footnote{Nota do editor: a passagem seguinte
foi desenvolvida.} e como
\(\SGARefLink{sga3:viii:corollary:4:7:eq:star:01}{(*)}\) mostra que
\(\operatorname{Spec}(A_0)\) é coberto pelos abertos afins
\(D(f_ig_i)\), reduzimo-nos ao caso em que há um elemento \(f\in A_1\)
invertível em \(A\). Então, para todo \(n\in\mathbb Z\), \(f^n\) é um
elemento invertível de \(A_n\), e define portanto um isomorfismo
\[
  A_0\xrightarrow{\ \sim\ }A_n,
  \qquad h\longmapsto f^nh.
\]
Além disso,
